\lysh{15987_SOLUTION}{1}
\hfill \par
ΛΥΣΗ
\par
α) Η ευθεία που διέρχεται από τα $A,B$ θα έχει εξίσωση: $y = \lambda x + b$ αφού $x_A \ne x_B$.
\[ \text{Οπότε } \lambda_{AB} = \frac{y_B - y_A}{x_B - x_A} = \frac{3 - 1}{2 - 1} = \frac{2}{1} = 2 \]
Άρα η ευθεία θα είναι της μορφής $y = 2x + b$. Αντικαθιστώντας τις συντεταγμένες του σημείου $A$ θα έχουμε: $1 = 2 \cdot 1 + b \Leftrightarrow b = -1$. Επομένως, η εξίσωση της ευθείας θα είναι η $(\varepsilon): y = 2x - 1$.

\par
β)
\par
α' τρόπος
\par
Το σημείο $E(3,5)$ ανήκει στην $(\varepsilon)$, διότι $2 \cdot 3 - 1 = 5$. Το σημείο $\Gamma$ βρίσκεται στην ίδια ευθεία παράλληλη στον άξονα $x'x$ την $y=5$ με το $E$ και «δεξιά» από αυτήν, ενώ το σημείο $O(0,0)$ βρίσκεται στο άλλο ημιεπίπεδο, «αριστερά» από αυτήν, όπως βλέπουμε και στο επόμενο σχήμα.
% TODO: TikZ σχήμα (μην το κατασκευάσεις)

\par
β' τρόπος
\par
Για $y=0$ το σημείο $A\left(\frac{1}{2}, 0\right)$ ανήκει στην ευθεία $(\varepsilon)$ και το διάνυσμα $\vec{AO} = \left(-\frac{1}{2}, 0\right)$ είναι παράλληλο στον άξονα $x'x$.
\par
Επίσης το σημείο $E(3,5)$ ανήκει στην $(\varepsilon)$, διότι $2 \cdot 3 - 1 = 5$ και το διάνυσμα $\vec{E\Gamma} = (2^{100} - 3, 5 - 5) = (2^{100} - 3, 0)$ είναι παράλληλο στον άξονα $x'x$. Τα διανύσματα $\vec{AO}, \vec{E\Gamma}$ έχουν αρχή στην ευθεία $(\varepsilon)$ και είναι αντίρροπα, οπότε ανήκουν σε διαφορετικά ημιεπίπεδα της $(\varepsilon)$. Συνεπώς και τα σημεία $O, \Gamma$ δεν ανήκουν στο ίδιο ημιεπίπεδο από αυτά που ορίζει η $(\varepsilon)$.

\par
γ) Τα τρίγωνα $AOB$ και $A\Gamma B$ έχουν την ίδια βάση $AB$, με φορέα την ευθεία $(\varepsilon)$.
\par
Η απόσταση του $O$ και του $\Gamma$ αντίστοιχα από την $(\varepsilon)$ είναι:
\[ d(O, \varepsilon) = \frac{|2 \cdot 0 - 0 - 1|}{\sqrt{2^2 + (-1)^2}} = \frac{1}{\sqrt{5}} \]
\[ d(\Gamma, \varepsilon) = \frac{|2 \cdot 2^{100} - 5 - 1|}{\sqrt{2^2 + (-1)^2}} = \frac{|2^{101} - 6|}{\sqrt{5}} = \frac{2^{101} - 6}{\sqrt{5}} \]
Συνεπώς $d(O, \varepsilon) < d(\Gamma, \varepsilon)$ άρα το τρίγωνο $A\Gamma B$ έχει μεγαλύτερο εμβαδόν από το $AOB$.
\par
\hfill \par
}